\lysh{21883_SOLUTION}{1}
\textbf{ΛΥΣΗ}

$\alpha)$ Η παράμετρος της παραβολής είναι $p=2$, άρα η εστία είναι το $E(0, \frac{p}{2}) = (0,1)$ και η εξίσωση της διευθετούσας $\delta: y=-1$.

Για τα κοινά σημεία της παραβολής $x^2=4y$ και της ευθείας $y=x-2$:
$x^2 = 4(x-2)$
$x^2 = 4x-8$
$x^2-4x+8=0$
Η διακρίνουσα είναι $\Delta = (-4)^2 - 4(1)(8) = 16 - 32 = -16 < 0$.
Εφόσον $\Delta < 0$, η εξίσωση δεν έχει πραγματικές λύσεις, άρα η ευθεία και η παραβολή δεν έχουν κοινά σημεία.

$\beta)$ Τα κοινά σημεία της παραβολής και της ευθείας είναι οι λύσεις του συστήματος των εξισώσεών τους.
% TODO: TikZ σχήμα

$\gamma)$
i. Έστω σημείο $M(x,y)$ σημείο της παραβολής, οπότε $M(x, \frac{1}{4}x^2)$. Η απόστασή του από την ευθεία $\varepsilon: x-y-2=0$ είναι
\[
d(M,\varepsilon) = \frac{|x - \frac{1}{4}x^2 - 2|}{\sqrt{1^2 + (-1)^2}} = \frac{|-( \frac{1}{4}x^2 - x + 2 )|}{\sqrt{2}} = \frac{|\frac{1}{4}x^2 - x + 2|}{\sqrt{2}}
\]
διότι η διακρίνουσα του τριωνύμου $\frac{1}{4}x^2 - x + 2$ είναι $\Delta = (-1)^2 - 4 \cdot \frac{1}{4} \cdot 2 = 1-2 = -1 < 0$ και άρα $\frac{1}{4}x^2 - x + 2 > 0$ για κάθε πραγματικό αριθμό $x$.

ii. Επομένως
\[
d(M,\varepsilon) = \frac{\frac{1}{4}x^2 - x + 2}{\sqrt{2}} = \frac{1}{\sqrt{2}} \left( \left(\frac{1}{2}x-1\right)^2 + 1 \right)
\]
Η απόσταση του M από την ευθεία γίνεται ελάχιστη όταν $\left( \frac{1}{2}x-1 \right)^2 = 0 \Leftrightarrow \frac{1}{2}x-1=0 \Leftrightarrow x=2$.

Άρα η ελάχιστη απόσταση $d(M,\varepsilon) = \frac{\sqrt{2}}{2}$ όταν $x=2$. Επομένως το ζητούμενο σημείο της παραβολής που απέχει ελάχιστη απόσταση από την ευθεία $\varepsilon$ είναι το $M(2,1)$.
\end{lysh}